\chapter{Suicidal thoughts}
\index{Suicidal thoughts}

\medskip
\noindent\textit{Educational reference; always seek professional advice for individual care.}

\section*{Definition and Overview}
Suicidal thoughts are thoughts about ending one's own life. They range enormously in intensity, from a fleeting, passive wish to not wake up, through a persistent feeling that life is not worth living, to active, detailed plans to die. Suicidal thinking exists on a spectrum, and it is far more common than most people realise. Having such thoughts does not mean a person will act on them, but it is always a signal that something is wrong and that the person is in real distress. The thoughts deserve to be taken seriously, spoken about openly, and responded to with care rather than judgement or alarm.

People who experience suicidal thoughts are often suffering from a heavy combination of pain, hopelessness, and the feeling that they are a burden. It is important to understand that these thoughts are a symptom of distress, not a sign of weakness, and that they can pass. With the right support, the great majority of people who have had suicidal thoughts do not die by suicide, and many recover fully. This chapter explains what suicidal thoughts are, why they occur, how they are recognised and assessed, and how the person, their family, and health professionals can respond. It is written to reassure people who are struggling that they are not alone, that help works, and that reaching out is a sign of strength.

\section*{Epidemiology}
Suicidal thoughts are common. International research consistently finds that roughly one in ten to one in twelve adults has had suicidal thoughts at some point in their life, and the figures are similar in many countries. They are more common than suicide attempts, which are in turn far more common than deaths by suicide, and each of these steps on the continuum is a chance to intervene. Suicidal thoughts can begin at any age, including in childhood, and they become more frequent through the teenage years. Young people, women, people who are socially isolated, and people living with mental illness, chronic pain, or significant stress are all at higher risk, and exposure to suicide, whether in the community, the media, or among family, can increase vulnerability.

The rate of death by suicide in many countries has remained broadly steady over recent decades, at around ten to twelve deaths per 100,000 people per year, and it remains the leading cause of death among young adults in some age groups. Men die by suicide at higher rates than women, while women attempt suicide more often. These patterns underline that suicidal thoughts are a public health problem that touches every community, and that early, compassionate responses at the level of the individual are the most important tool we have. Most people who have suicidal thoughts do eventually find relief, which is the most hopeful statistic of all.

\section*{Aetiology and Pathophysiology}
Suicidal thoughts arise from a combination of factors rather than a single cause. The strongest drivers are usually mental health conditions, particularly depression, but anxiety, post-traumatic stress, psychosis, and substance misuse also play major roles. The thoughts are commonly amplified by life circumstances: the end of a relationship, bereavement, financial crisis, job loss, bullying, or chronic health problems. A person's sense of hopelessness, the belief that things cannot improve, and the feeling of being a burden to others are psychological factors that distinguish serious suicidal thinking from the ordinary low moments everyone experiences.

There is also a biological dimension. Suicide risk runs in families, and research points to inherited variations in the way the brain processes stress, impulse control, and mood. Stress hormones, the serotonin system, and the brain circuits involved in decision-making and impulse control are all implicated. Importantly, vulnerability fluctuates: a person may carry risk factors for years but become acutely suicidal only during a period of crisis, illness, intoxication, or a specific triggering event. This is why asking directly about suicidal thoughts at the right moment can be life-saving, and why a period of intense distress, even in someone who has coped well for years, should never be dismissed.

\begin{itemize}
    \item \textbf{Depression:} the single most important risk factor; persistent low mood, hopelessness, and self-criticism fuel suicidal thinking
    \item \textbf{Other mental health conditions:} bipolar disorder, schizophrenia, anxiety, post-traumatic stress disorder, and eating disorders all raise risk
    \item \textbf{Substance misuse:} alcohol and drugs impair judgement, worsen mood, and are involved in a large share of suicide attempts
    \item \textbf{Trauma and abuse:} childhood abuse, assault, and violence in adult relationships increase vulnerability
    \item \textbf{Bereavement and loss:} grief, particularly after suicide, is a high-risk period
    \item \textbf{Chronic illness and pain:} physical illness, chronic pain, and disability are associated with suicidal thoughts
    \item \textbf{Life stress:} relationship breakdown, financial problems, legal trouble, bullying, and social isolation
    \item \textbf{Previous attempts:} a history of self-harm or a previous suicide attempt is one of the strongest predictors of future risk
\end{itemize}

\section*{Clinical Features}
The "features" of suicidal thoughts are the signs a person may be experiencing them, whether they say so or not.

\begin{itemize}
    \item \textbf{Direct statements:} talking about dying, killing oneself, or wanting it to end, or saying goodbye to people, which must always be taken seriously
    \item \textbf{Indirect expressions:} comments such as "you won't have to worry about me for much longer"
    \item \textbf{Hopelessness:} a persistent belief that things cannot change or improve
    \item \textbf{Feeling trapped or a burden:} believing there is no way out, or that loved ones would be better off without them
    \item \textbf{Withdrawal:} pulling away from friends, family, and activities once enjoyed
    \item \textbf{Changes in mood:} deep sadness, irritability, or anxiety, and in some cases a sudden calm that can signal a decision has been made
    \item \textbf{Neglect of self:} poor sleep, eating changes, giving away possessions, or losing interest in the future
    \item \textbf{Increased risk-taking:} dangerous driving, reckless substance use, or self-harm
    \item \textbf{Searching for means:} obtaining medicines, weapons, or researching ways to die, which raises risk substantially
\end{itemize}

\section*{Differential Diagnosis}
\begin{itemize}
    \item \textbf{Depression:} the commonest context; low mood, anhedonia, and self-worth issues underpin most suicidal thinking
    \item \textbf{Anxiety disorders:} severe anxiety, panic, and agitation can be accompanied by desperate thoughts of escape
    \item \textbf{Post-traumatic stress disorder:} intrusive memories and re-experiencing can drive hopelessness and suicidal thoughts
    \item \textbf{Psychosis:} voices commanding self-harm or persecutory beliefs can be a cause, requiring urgent specialist care
    \item \textbf{Substance intoxication and withdrawal:} alcohol and drugs can produce or amplify suicidal thoughts, and withdrawal from alcohol or sedatives is a high-risk period
    \item \textbf{Personality and developmental factors:} impulsivity and borderline personality traits can present with crisis-level distress
    \item \textbf{Medical conditions:} thyroid disease, chronic pain, and neurological conditions can cause or worsen mood disturbance and suicidal thinking
\end{itemize}

\section*{When to Seek Medical Care}
If you are having suicidal thoughts, or if you are worried about someone who is, the most important step is to talk to someone you trust and to seek professional help. A general practitioner is a good first stop, and speaking to a crisis counsellor is immediate and free. Asking someone directly whether they are having suicidal thoughts does not put the idea in their head; it can be a relief and can open the door to help.

If you need support right now, contact a trusted person, a local crisis line, an urgent mental-health service, or the nearest hospital emergency department. In an emergency, contact your local emergency services. Crisis-line availability and telephone numbers vary by country.

\textbf{Contact your local emergency services for:}
\begin{itemize}
    \item Thoughts of suicide together with a specific plan and access to the means to act
    \item Taking an overdose, self-harm that requires medical attention, or any attempt to end life
    \item A statement that the person is about to end their life, or behaviour suggesting they will
    \item Someone who is out of contact after expressing suicidal thoughts
\end{itemize}

\section*{Diagnosis and Investigations}
\begin{itemize}
    \item \textbf{Direct, gentle questioning:} asking about thoughts of self-harm and suicide openly and without judgement is the cornerstone of assessment; it is not intrusive when done with care
    \item \textbf{Detailed history:} the nature, frequency, and intensity of the thoughts, any plans, access to means, and previous attempts are explored
    \item \textbf{Mental state examination:} mood, thinking, hopelessness, agitation, and insight are assessed
    \item \textbf{Assessment of protective factors:} reasons for living, supportive relationships, and future plans are important parts of the picture
    \item \textbf{Risk and safety planning:} a written plan agreed with the person, covering warning signs, coping strategies, supports, and when to seek urgent help
    \item \textbf{Physical review:} blood tests and examination are used when substance use, medication, or physical illness may be contributing
    \item \textbf{Specialist referral:} psychiatrists, psychologists, and public mental health services provide comprehensive assessment and treatment
\end{itemize}

\section*{Management}
\begin{itemize}
    \item \textbf{Talking openly:} being listened to without judgement is itself therapeutic, and family and friends can provide this support
    \item \textbf{Psychological therapy:} cognitive behaviour therapy, and specifically therapies designed for suicidal thinking, help people find alternatives to hopelessness and build coping skills
    \item \textbf{Treatment of underlying conditions:} effective treatment of depression, anxiety, trauma, or substance misuse is the most important step
    \item \textbf{Medication:} antidepressants and other medicines are effective for many people, though some take weeks to work and some can temporarily increase risk in young people, which is why careful supervision matters
    \item \textbf{Safety planning:} a practical, personalised plan that limits access to means, identifies warning signs, and lists who to call
    \item \textbf{Reducing access to means:} safe storage of medicines, firearms, and other means saves lives by creating time and distance between impulse and action
    \item \textbf{Hospital care:} people at acute risk may be admitted briefly so they are safe while treatment begins
    \item \textbf{Peer and community support:} support groups and services designed for specific groups reduce isolation
    \item \textbf{Ongoing follow-up:} regular contact with a GP or mental health professional, because risk rises and falls and support needs to be available over time
\end{itemize}

\section*{Complications}
\begin{itemize}
    \item \textbf{Suicide attempt:} the most serious consequence; attempts can cause lasting physical injury and, in some cases, death
    \item \textbf{Death by suicide:} the outcome everyone hopes to prevent, and one that the combination of care and safety measures can reduce
    \item \textbf{Injury and disability:} surviving an attempt with severe physical injury can create new, long-term challenges
    \item \textbf{Worsening mental health:} unaddressed suicidal thinking can deepen depression and increase the risk of future attempts
    \item \textbf{Impact on family and friends:} those close to a person in crisis experience significant distress, and bereavement after suicide is especially complicated and painful
    \item \textbf{Repeated crises:} without treatment and support, suicidal thoughts tend to recur during future periods of stress
\end{itemize}

\section*{Prognosis and Outcomes}
The outlook for a person with suicidal thoughts is very good when the thoughts are taken seriously and the person is connected to help. Suicidal thinking is overwhelmingly a temporary state; it waxes and wanes, and it responds to treatment of the underlying condition and to support during crisis. Studies consistently show that people who reach out for help, engage in therapy, and have the support of a network recover, and the large majority of people who have had suicidal thoughts never go on to die by suicide. Treatment works but takes time, and early improvements can be followed by a fragile period before the effect is complete, which is why close follow-up and safety planning are built into good care. Recovery is possible, and it happens every day.

\section*{Prevention and Self-Care}
\begin{itemize}
    \item Reach out early; talking to a GP, a counsellor, or a trusted person is the single most protective step
    \item If you are caring for someone who is struggling, ask directly about suicidal thoughts, listen without judgement, and take what they say seriously
    \item Keep local crisis-line and emergency numbers saved in your phone, and make sure someone you trust knows how to reach you
    \item Remove or safely store means of self-harm, including medicines and firearms, especially during a crisis
    \item Protect sleep, maintain basic routines, and limit alcohol and drugs, which make thinking worse
    \item Stay connected; isolation deepens hopelessness, and a small number of safe, regular contacts make a difference
    \item Make a written safety plan with your doctor and keep it where you can see it
    \item Be patient and kind with yourself; recovery from suicidal thoughts is a process, and setbacks do not erase progress
\end{itemize}

\section*{Sources and Further Reading}
The following organisations provide reliable information, support, and crisis services for people experiencing suicidal thoughts.

\begin{itemize}
    \item World Health Organization. \textit{Suicide prevention.} https://www.who.int/health-topics/suicide
    \item World Health Organization. \textit{Preventing suicide: a resource for media professionals and others working with them.} https://www.who.int/publications/i/item/9789240076846
    \item Use the crisis-line directory or emergency-service information provided by your local health authority.
\end{itemize}

\clearpage
